\RequirePackage{fix-cm}
\documentclass[11pt]{article}
\usepackage[margin=1in]{geometry}
\usepackage{fancyhdr}
\usepackage{enumitem}
\usepackage{parskip}
\usepackage{graphicx}
\usepackage{needspace}

\pagestyle{fancy}
\fancyhf{}
\fancyhead[C]{\textbf{SEG2105}}
\fancyfoot[C]{\thepage}
\setlength{\headheight}{14pt}
\setlength{\emergencystretch}{2em}

\newcounter{chapterNumber}
\setcounter{chapterNumber}{1}
\newcounter{headerOneCount}
\newcounter{headerTwoCount}[headerOneCount]
\setcounter{headerOneCount}{-1}
\renewcommand{\theheaderOneCount}{\arabic{headerOneCount}}
\renewcommand{\theheaderTwoCount}{\arabic{headerTwoCount}}

\newcommand{\Title}[1]{%
  \begin{center}
    \fontsize{16}{28}\selectfont\textbf{#1}
  \end{center}
}

\newcommand{\HeaderOne}[1]{%
  \Needspace{5\baselineskip}
  \refstepcounter{headerOneCount}%
  \setcounter{headerTwoCount}{-1}%
  {\noindent\fontsize{14}{22}\selectfont\textbf{\arabic{chapterNumber}.\theheaderOneCount\ #1}\par}
}

\newcommand{\HeaderTwo}[1]{%
  \Needspace{5\baselineskip}
  \refstepcounter{headerTwoCount}%
  {\noindent\fontsize{14}{20}\selectfont\textbf{\arabic{chapterNumber}.\theheaderOneCount.\theheaderTwoCount\ #1}\par}
}

\newcommand{\BodyText}[1]{%
  \fontsize{11}{15}\selectfont #1\par

}

\newenvironment{NumberedList}{%
  \begin{enumerate}[leftmargin=2em, itemsep=0.35em, topsep=0.35em]
}{%
  \end{enumerate}
}

\newenvironment{DottedList}{%
  \begin{itemize}[leftmargin=2em, itemsep=0.35em, topsep=0.35em, label=\textbullet]
}{%
  \end{itemize}
}

\begin{document}
\Title{Chapter 1: Software and Software Engineering}
\HeaderOne{The Nature of Software}
\begin{DottedList}
  \item Software is intangible: effort, progress, and quality can be difficult to see.
  \item Reproduction is inexpensive compared with development. Development is labour-intensive and hard to automate completely.
  \item Easy modification can hide expensive consequences: changes made without understanding the design can introduce defects and complexity.
  \item Software does not physically wear out. Its structure can deteriorate through poorly planned changes.
\end{DottedList}
\HeaderTwo{Types of software}
\begin{DottedList}
  \item Real-time embedded systems monitor or control equipment and must meet timing requirements; safety may be critical.
  \item Data-processing systems support business operations, where data accuracy and security matter.
  \item Games, mobile applications, and web applications have different users, operating environments, and quality priorities.
\end{DottedList}
\HeaderOne{What is Software Engineering?}
\BodyText{Software engineering solves customer problems through systematic development and evolution of large, sufficiently high-quality software systems within cost, time, and other constraints.}
\begin{DottedList}
  \item Understand the customer's problem before choosing a solution. Buying or adapting existing software may be appropriate.
  \item Apply established techniques in an organized way. Much development evolves existing systems.
  \item Divide work across a team and coordinate the interfaces between parts. One person cannot understand every detail of a large system.
  \item Estimate resources realistically and ensure benefits justify costs. Extra features can increase complexity without solving the problem.
\end{DottedList}
\HeaderOne{The Software Engineering Profession}
\BodyText{The term software engineering dates to 1968. An engineering approach combines technical knowledge, disciplined practice, and responsibility for the resulting systems.}
\HeaderTwo{Ethics}
\begin{DottedList}
  \item Protect the public interest and act responsibly toward clients and employers.
  \item Aim for high product and maintenance standards, and preserve independent professional judgment.
  \item Promote ethical management and the integrity of the profession.
  \item Support colleagues fairly and continue learning throughout a career.
\end{DottedList}
\HeaderOne{Software Quality}
\begin{DottedList}
  \item Usability: users can learn the system and complete their work effectively.
  \item Efficiency: the system avoids unnecessary use of processor time, memory, and other resources.
  \item Reliability: the system performs required functions without failure under stated conditions.
  \item Maintainability: developers can understand and change the system without excessive effort.
  \item Reusability: components can serve other projects or applications.
\end{DottedList}
\HeaderTwo{Stakeholders and tradeoffs}
\begin{DottedList}
  \item Users value ease of learning and effective task completion; customers value solving a problem at an acceptable total cost.
  \item Developers value understandable designs and components that are easy to maintain and reuse.
  \item Development managers balance customer satisfaction with development and maintenance costs.
  \item Quality goals can conflict. Optimization may make code harder to maintain; a richer interface may consume more resources.
  \item Set quality objectives before designing. Over-engineering beyond those needs wastes resources.
\end{DottedList}
\HeaderOne{Software Engineering Projects}
\begin{DottedList}
  \item Corrective maintenance fixes defects.
  \item Adaptive maintenance responds to environmental changes such as operating systems, databases, or external rules.
  \item Enhancement adds user-visible capabilities.
  \item Reengineering or perfective work improves internal structure and maintainability.
  \item Greenfield development creates a new system. Many projects instead evolve legacy systems.
\end{DottedList}
\HeaderOne{Activities Common to Software Projects}
\HeaderTwo{Requirements and specification}
\begin{DottedList}
  \item Domain analysis develops understanding of the field and problem.
  \item Gathering obtains information from stakeholders and other sources.
  \item Analysis organizes and resolves that information. Specification records required behaviour and constraints.
\end{DottedList}
\HeaderTwo{Design}
\begin{DottedList}
  \item Systems engineering allocates responsibilities between hardware and software.
  \item Architecture divides software into subsystems and defines their interactions.
  \item Detailed design determines subsystem internals, user interfaces, and database structures.
\end{DottedList}
\HeaderTwo{Modelling, implementation, and delivery}
\begin{DottedList}
  \item Use-case models describe goals and interactions. Structural models describe organization. Behavioural models describe interactions and changes over time.
  \item Programming implements the design. Reviews, inspections, and tests assess quality.
  \item Deployment makes the system available. Process management coordinates work, resources, and progress.
  \item These activities recur and inform one another as software evolves.
\end{DottedList}
\end{document}
